\section{Conclusión}

La documentación de la arquitectura técnica del sistema permite comprender que las 
decisiones de diseño adoptadas no son arbitrarias, sino que responden a principios 
consistentes de separación de responsabilidades que se aplican de forma simétrica en 
ambas capas del sistema.

En el backend, la división entre DTOs, repositorios y controladores garantiza que 
cada componente tenga una responsabilidad acotada y bien definida: el DTO transforma, 
el repositorio accede a los datos y el controlador orquesta. Esta uniformidad hace 
que cada módulo sea predecible y que añadir un nuevo servicio al sistema implique 
replicar un patrón ya establecido, sin necesidad de tomar decisiones arquitectónicas 
nuevas en cada iteración.

En el frontend, la misma lógica se expresa a través de la separación entre 
\texttt{types/}, \texttt{services/} y las stores de Pinia: los tipos definen el 
contrato, los servicios ejecutan la comunicación y las stores mantienen el estado. 
Esta cadena garantiza que los componentes Vue puedan consumir datos de la API de 
forma tipada y segura, sin acoplarse a los detalles del protocolo HTTP ni a la 
convención de nomenclatura del backend.

El flujo de consumo de la API sintetiza estas decisiones en una secuencia coherente 
que atraviesa todas las capas del sistema. Comprender este flujo es fundamental para 
interpretar correctamente el código implementado en los hitos siguientes, ya que cada 
módulo desarrollado, tanto en backend como en frontend, sigue exactamente este mismo 
patrón de comunicación.